% Paper I: Fejér–heat generators and Lipschitz control for the Weil quadratic functional
% Style: Rodgers–Tao (Forum of Mathematics, Pi)
\documentclass[11pt,a4paper]{article}

\usepackage{amsmath,amssymb,amsthm}
\usepackage{mathrsfs}
\usepackage{hyperref}
\usepackage{enumitem}
\usepackage{booktabs}
\usepackage{geometry}
\geometry{margin=1in}

% Theorem environments (Tao style)
\newtheorem{theorem}{Theorem}[section]
\newtheorem{lemma}[theorem]{Lemma}
\newtheorem{proposition}[theorem]{Proposition}
\newtheorem{corollary}[theorem]{Corollary}
\theoremstyle{definition}
\newtheorem{definition}[theorem]{Definition}
\theoremstyle{remark}
\newtheorem{remark}[theorem]{Remark}

% Custom commands
\newcommand{\R}{\mathbb{R}}
\newcommand{\C}{\mathbb{C}}
\newcommand{\Z}{\mathbb{Z}}
\newcommand{\N}{\mathbb{N}}
\DeclareMathOperator{\supp}{supp}
\DeclareMathOperator{\sgn}{sgn}

\title{Fej\'er--heat generators and Lipschitz control\\for the Weil quadratic functional}

\author{
  {\sc First Author}\thanks{Department of Mathematics, University A; email: author1@university.edu}
  \and
  {\sc Second Author}\thanks{Department of Mathematics, University B; email: author2@university.edu}
}

\date{}

\begin{document}

\maketitle

\begin{abstract}
For a test function $\Phi$ in a suitable cone of even, compactly supported functions, the Weil quadratic functional $Q(\Phi)$ arises naturally from the Guinand--Weil explicit formula in analytic number theory. We establish three foundational results for the operator-theoretic study of this functional. First, we fix a normalization convention (T0) that identifies the Guinand--Weil frequency axis with the repository convention $\xi = \eta/(2\pi)$, ensuring compatibility with classical references. Second, we prove a cone density theorem (A1'): on any frequency compact $[-K, K]$, the cone generated by Fej\'er--heat windows $\Phi_{B,t}(\xi) = (1 - |\xi|/B)_+ \cdot e^{-4\pi^2 t \xi^2}$ is dense in the cone of continuous, even, nonnegative functions under the uniform norm. The approximation is explicit, with bandwidth $B = O(1/\varepsilon)$ and heat scale $t$ depending on the mollification error. Third, we establish Lipschitz continuity (A2) of $Q$ on each compact window: for $\Phi_1, \Phi_2 \in C_c([-K,K])$ even, one has $|Q(\Phi_1) - Q(\Phi_2)| \leq L_Q(K) \|\Phi_1 - \Phi_2\|_\infty$, where the Lipschitz constant $L_Q(K)$ is computed explicitly from the Archimedean density $a(\xi)$ and the finitely many prime weights in $[-K,K]$. The Fej\'er--heat tail outside the support contributes at most $t^{-1} e^{-t(\log N)^2}$, which is absorbed into the continuity budget. These results form the analytic foundation for the Toeplitz barrier and RKHS contraction arguments developed in [Part~II]. This paper is the first of three developing an operator-theoretic approach to the Weil positivity criterion.
\end{abstract}

\tableofcontents

%=============================================================================
\section{Introduction}
%=============================================================================

\subsection{The Weil functional and prime distribution}

Let $\Lambda$ denote the von Mangoldt function, and for $n \geq 2$ define the prime nodes and Weil weights
\begin{equation}\label{eq:nodes-weights}
  \xi_n := \frac{\log n}{2\pi}, \qquad w_Q(n) := \frac{2\Lambda(n)}{\sqrt{n}}.
\end{equation}
For an even test function $\Phi : \R \to \R$ supported on a compact interval $[-K, K]$, the \emph{Weil quadratic functional} is
\begin{equation}\label{eq:Q-def}
  Q(\Phi) := \int_{-K}^{K} a^*(\xi)\, \Phi(\xi)\, d\xi - \sum_{\xi_n \in [-K,K]} w_Q(n)\, \Phi(\xi_n),
\end{equation}
where
\begin{equation}\label{eq:a-star}
  a^*(\xi) := 2\pi \left( \log \pi - \Re \psi\left(\tfrac{1}{4} + i\pi\xi\right) \right)
\end{equation}
is the \emph{Archimedean density} arising from the Guinand--Weil explicit formula~\cite{Weil1952,Guinand1948}, and $\psi(z) = \Gamma'(z)/\Gamma(z)$ denotes the digamma function. The functional $Q$ encodes information about the distribution of prime numbers through the Fourier-analytic structure of the explicit formula. Understanding when $Q$ is nonnegative on natural test cones is a central problem in analytic number theory.

Our approach to this problem decomposes into a sequence of analytic modules:

\begin{table}[h]
\centering
\caption{Analytic modules for the Weil positivity analysis.}
\label{tab:modules}
\begin{tabular}{@{}lll@{}}
\toprule
Module & Statement & Reference \\
\midrule
(T0)  & Normalization crosswalk: Guinand--Weil $\leftrightarrow$ repository & This paper \\
(A1') & Fej\'er--heat cone dense in $C^+_{\mathrm{even}}([-K,K])$ & This paper \\
(A2)  & $Q$ Lipschitz on $C_c([-K,K])$ & This paper \\
(A3)  & Toeplitz barrier: $\inf_\theta P_A(\theta) \geq c_* > 0$ & [Part~II] \\
(RKHS) & Prime cap: $\|T_P\| \leq \rho(t) < c_*/4$ & [Part~II] \\
\bottomrule
\end{tabular}
\end{table}

\noindent The synthesis of these modules appears in [Part~III].

\subsection{Main results}

We now state the main results of this paper.

\begin{theorem}[T0: Normalization crosswalk]\label{thm:T0}
Under the Fourier convention $\widehat{\varphi}(\xi) = \int_\R \varphi(t) e^{-2\pi i t \xi}\, dt$, the repository quadratic functional $Q(\varphi)$ coincides with the classical Guinand--Weil functional $Q_{\mathrm{GW}}(\varphi_{\mathrm{GW}})$ under the change of variables $\eta = 2\pi\xi$, $\varphi_{\mathrm{GW}}(\eta) = \varphi(\eta/(2\pi))$. The Archimedean density transforms as $a^*(\xi) = 2\pi\, a(\xi)$ where $a(\xi) = \log\pi - \Re\psi(\tfrac{1}{4} + i\pi\xi)$.
\end{theorem}

The test functions we consider are \emph{Fej\'er--heat windows}:
\begin{equation}\label{eq:Phi-def}
  \Phi_{B,t}(\xi) := \left(1 - \frac{|\xi|}{B}\right)_+ \cdot e^{-4\pi^2 t \xi^2},
\end{equation}
where $(x)_+ := \max(x, 0)$, $B > 0$ is the bandwidth, and $t > 0$ is the heat scale. For $\tau \in [-K, K]$, the symmetrized atoms are
\begin{equation}\label{eq:atom}
  \Phi_{B,t,\tau}(\xi) := \frac{1}{2}\left[ \Lambda_B(\xi - \tau) \rho_t(\xi - \tau) + \Lambda_B(\xi + \tau) \rho_t(\xi + \tau) \right],
\end{equation}
where $\Lambda_B(x) = (1 - |x|/B)_+$ is the Fej\'er kernel (tent function) and $\rho_t(x) = (4\pi t)^{-1/2} e^{-x^2/(4t)}$ is the heat kernel.

\begin{theorem}[A1': Cone density]\label{thm:A1-density}
Let $K = [-R, R]$ with $R > 0$. The closed convex cone $\mathcal{C}$ generated by finite nonnegative combinations of the symmetrized atoms $\{\Phi_{B,t,\tau}\}$ with $\tau \in [-R, R]$ and $B \geq B_0(R, \varepsilon)$ is dense in $C^+_{\mathrm{even}}([-R, R])$ under the uniform norm $\|\cdot\|_\infty$.
\end{theorem}

\noindent The explicit dependence is $B_0 = O(R/\varepsilon)$; see Lemma~\ref{lem:fejer-approx} for the precise constant.

\begin{theorem}[A2: Lipschitz continuity]\label{thm:A2-Lip}
Fix $K = [-R, R]$. For even $\Phi_1, \Phi_2 \in C_c([-R, R])$, one has
\begin{equation}\label{eq:Lip-bound}
  |Q(\Phi_1) - Q(\Phi_2)| \leq L_Q(K)\, \|\Phi_1 - \Phi_2\|_\infty,
\end{equation}
where the Lipschitz constant is
\begin{equation}\label{eq:LQ}
  L_Q(K) := \|a^*\|_{L^1(K)} + \sum_{\xi_n \in K} w_Q(n).
\end{equation}
The sum is finite because only $n \leq e^{2\pi R}$ satisfy $\xi_n \in K$.
\end{theorem}

\subsection{Proof overview}

The proof of Theorem~\ref{thm:T0} is a direct calculation: the Jacobian $d\eta = 2\pi\, d\xi$ absorbs into the density $a^*$, and the evenness of $\varphi$ ensures the prime sum matches under $\xi_n = (\log n)/(2\pi)$. We include it for completeness and to fix notation used throughout the trilogy.

The density theorem (A1') proceeds in three steps:
\begin{enumerate}[label=(\roman*)]
  \item \emph{Mollification}: Convolve with the heat kernel $\rho_t$ to produce $g = \widetilde{f} * \rho_t \in C^\infty$, with $\|g - f\|_\infty < \varepsilon/3$ for small $t$.
  \item \emph{Positive Riemann sums}: Approximate $g$ by a finite sum of heat translates with nonnegative coefficients $g(\tau_j^*)$.
  \item \emph{Fej\'er truncation}: Replace the unbounded Gaussian by the compactly supported Fej\'er kernel $\Lambda_B$, with $B$ large enough that $|\Lambda_B(\xi - \tau) - 1| < \varepsilon/3$ on $[-R, R]$.
\end{enumerate}
The triangle inequality collects the errors. The symmetrization preserves evenness.

The Lipschitz bound (A2) follows from two observations: (i)~the Archimedean density $a^*$ is bounded on $[-R, R]$, hence the integral term is Lipschitz in $\Phi$; (ii)~the prime sum is a finite sum of point evaluations, and finiteness is guaranteed by $\xi_n \in [-R, R] \Leftrightarrow n \leq e^{2\pi R}$.

\subsection{Notation}\label{sec:notation}

We collect the notational conventions used throughout.

\begin{itemize}[leftmargin=2em]
  \item \textbf{Prime nodes and weights}: $\xi_n = (\log n)/(2\pi)$, $w_Q(n) = 2\Lambda(n)/\sqrt{n}$.
  \item \textbf{Bandwidth window}: $W_K = [-K, K]$, the frequency compact.
  \item \textbf{Fej\'er kernel}: $\Lambda_B(\xi) = (1 - |\xi|/B)_+$, supported on $[-B, B]$.
  \item \textbf{Heat kernel}: $\rho_t(x) = (4\pi t)^{-1/2} e^{-x^2/(4t)}$, with $\int_\R \rho_t = 1$.
  \item \textbf{Archimedean density}: $a(\xi) = \log\pi - \Re\psi(\tfrac{1}{4} + i\pi\xi)$, $a^*(\xi) = 2\pi\, a(\xi)$.
  \item \textbf{Asymptotic notation}: $O(\cdot)$ denotes absolute constants; $O_K(\cdot)$ allows dependence on $K$.
\end{itemize}

%=============================================================================
\section{Normalization: the Guinand--Weil crosswalk (T0)}
%=============================================================================

\subsection{Fourier conventions}

We adopt the Fourier normalization
\begin{equation}\label{eq:Fourier}
  \widehat{\varphi}(\xi) = \int_\R \varphi(t)\, e^{-2\pi i t \xi}\, dt, \qquad
  \varphi(t) = \int_\R \widehat{\varphi}(\xi)\, e^{2\pi i t \xi}\, d\xi,
\end{equation}
using the Lebesgue measure $d\xi$ on the frequency side. All test functions are even, so the sine components vanish.

\subsection{The Guinand--Weil functional}

The classical Guinand--Weil functional~\cite{Guinand1948,Weil1952} is defined on even test functions $\varphi_{\mathrm{GW}}$ on the frequency axis $\eta \in \R$ by
\begin{equation}\label{eq:Q-GW}
  Q_{\mathrm{GW}}(\varphi_{\mathrm{GW}}) := \int_\R \left( \log\pi - \Re\psi\left(\tfrac{1}{4} + \tfrac{i\eta}{2}\right) \right) \varphi_{\mathrm{GW}}(\eta)\, d\eta
  - \sum_{n \geq 2} \frac{\Lambda(n)}{\sqrt{n}}\, \bigl[ \varphi_{\mathrm{GW}}(\log n) + \varphi_{\mathrm{GW}}(-\log n) \bigr].
\end{equation}

\begin{proof}[Proof of Theorem~\ref{thm:T0}]
Set $\eta = 2\pi\xi$ and $\varphi(\xi) = \varphi_{\mathrm{GW}}(2\pi\xi)$. For the Archimedean integral:
\[
  \int_\R \left( \log\pi - \Re\psi\left(\tfrac{1}{4} + \tfrac{i\eta}{2}\right) \right) \varphi_{\mathrm{GW}}(\eta)\, d\eta
  = \int_\R 2\pi \left( \log\pi - \Re\psi\left(\tfrac{1}{4} + i\pi\xi\right) \right) \varphi(\xi)\, d\xi,
\]
using $d\eta = 2\pi\, d\xi$ and $\psi(\tfrac{1}{4} + \tfrac{i\eta}{2}) = \psi(\tfrac{1}{4} + i\pi\xi)$.

For the prime sum, $\varphi_{\mathrm{GW}}(\pm \log n) = \varphi(\pm \xi_n)$ with $\xi_n = (\log n)/(2\pi)$. Since $\varphi$ is even:
\[
  \varphi_{\mathrm{GW}}(\log n) + \varphi_{\mathrm{GW}}(-\log n) = \varphi(\xi_n) + \varphi(-\xi_n) = 2\varphi(\xi_n).
\]
Thus
\[
  \sum_{n \geq 2} \frac{\Lambda(n)}{\sqrt{n}}\, \bigl[ \varphi_{\mathrm{GW}}(\log n) + \varphi_{\mathrm{GW}}(-\log n) \bigr]
  = \sum_{n \geq 2} \frac{2\Lambda(n)}{\sqrt{n}}\, \varphi(\xi_n) = \sum_{n \geq 2} w_Q(n)\, \varphi(\xi_n).
\]
Combining, $Q(\varphi) = Q_{\mathrm{GW}}(\varphi_{\mathrm{GW}})$ as claimed.
\end{proof}

\begin{remark}
The doubling convention $w_Q(n) = 2\Lambda(n)/\sqrt{n}$ for nodes at $\xi_n > 0$ is equivalent to placing unit weights at both $\pm \xi_n$; evenness makes the two conventions identical.
\end{remark}

%=============================================================================
\section{Cone density on frequency compacts (A1')}
%=============================================================================

We establish that the Fej\'er--heat cone is dense in the cone of continuous, even, nonnegative functions on any compact.

\subsection{Heat mollification}

The heat kernel $\rho_t(x) = (4\pi t)^{-1/2} e^{-x^2/(4t)}$ satisfies:
\begin{enumerate}[label=(\alph*)]
  \item $\rho_t \geq 0$ and $\int_\R \rho_t = 1$ (probability density);
  \item $\rho_t \to \delta_0$ as $t \to 0^+$ in the sense of distributions;
  \item For $f \in C_b(\R)$, $\|f * \rho_t - f\|_\infty \to 0$ as $t \to 0^+$ (uniform approximation).
\end{enumerate}

\begin{lemma}[Mollification error]\label{lem:mollification}
Let $f \in C_c(\R)$ with $\supp f \subseteq [-L, L]$. For any $\varepsilon > 0$, there exists $t_0 = t_0(f, \varepsilon) > 0$ such that for all $t \in (0, t_0]$:
\[
  \|f * \rho_t - f\|_\infty < \varepsilon.
\]
Moreover, if $f \geq 0$ and $f$ is even, then $f * \rho_t \geq 0$ and $f * \rho_t$ is even.
\end{lemma}

\begin{proof}
Standard properties of approximate identities. Nonnegativity follows from $f \geq 0$ and $\rho_t \geq 0$. Evenness: if $f(x) = f(-x)$ and $\rho_t(x) = \rho_t(-x)$, then $(f * \rho_t)(-x) = (f * \rho_t)(x)$.
\end{proof}

\subsection{Fej\'er truncation}

\begin{lemma}[Fej\'er approximation]\label{lem:fejer-approx}
For $|\xi|, |\tau| \leq R$ and $B > 0$:
\[
  |\Lambda_B(\xi - \tau) - 1| \leq \frac{|\xi - \tau|}{B} \leq \frac{2R}{B}.
\]
In particular, choosing $B \geq 6R/\varepsilon$ ensures $|\Lambda_B(\xi - \tau) - 1| < \varepsilon/3$ on $[-R, R]^2$.
\end{lemma}

\begin{proof}
For $|x| \leq B$, $\Lambda_B(x) = 1 - |x|/B$, so $|\Lambda_B(x) - 1| = |x|/B$. Since $|\xi - \tau| \leq |\xi| + |\tau| \leq 2R$, the bound follows.
\end{proof}

\subsection{Proof of cone density}

\begin{proof}[Proof of Theorem~\ref{thm:A1-density}]
Fix $f \in C^+_{\mathrm{even}}([-R, R])$ and $\varepsilon > 0$. Extend $f$ by zero to $\widetilde{f} \in C_c(\R)$.

\textbf{Step 1 (Mollification).} By Lemma~\ref{lem:mollification}, there exists $t > 0$ such that $g := \widetilde{f} * \rho_t$ satisfies
\[
  \sup_{|\xi| \leq R} |g(\xi) - f(\xi)| < \varepsilon/3.
\]
The function $g$ is nonnegative, even, and smooth.

\textbf{Step 2 (Riemann sums).} Choose a uniform partition $-R = \tau_0 < \tau_1 < \cdots < \tau_N = R$ with mesh $\Delta$ small enough that
\[
  g_R(\xi) := \sum_{j=0}^{N-1} g(\tau_j^*)\, \rho_t(\xi - \tau_j^*)\, \Delta
\]
satisfies $\|g_R - g\|_{L^\infty([-R,R])} < \varepsilon/3$. The coefficients $g(\tau_j^*) \cdot \Delta$ are nonnegative.

\textbf{Step 3 (Fej\'er truncation).} Choose $B \geq 6R/\varepsilon$ per Lemma~\ref{lem:fejer-approx}. Define
\[
  h(\xi) := \sum_{j=0}^{N-1} g(\tau_j^*)\, \Delta\, \Phi_{B,t,\tau_j^*}(\xi).
\]
For $|\xi| \leq R$, $|\Lambda_B(\xi \pm \tau_j^*) - 1| < \varepsilon/3$, so
\[
  |h(\xi) - g_R^{\mathrm{sym}}(\xi)| \leq \frac{\varepsilon}{3} \sum_j g(\tau_j^*)\, \Delta \leq C_f \cdot \frac{\varepsilon}{3},
\]
where $C_f = \int_{-R}^R g(\tau)\, d\tau$. Rescaling $\varepsilon$ by $3\max(1, C_f)$:
\[
  \|h - g_R^{\mathrm{sym}}\|_\infty < \varepsilon/3.
\]

\textbf{Step 4 (Triangle inequality).} Since $g$ is even, $g_R^{\mathrm{sym}}$ satisfies $\|g_R^{\mathrm{sym}} - g\|_\infty < \varepsilon/3$. Combining:
\[
  \|h - f\|_\infty \leq \|h - g_R^{\mathrm{sym}}\|_\infty + \|g_R^{\mathrm{sym}} - g\|_\infty + \|g - f\|_\infty < \varepsilon.
\]
Since $h$ is a finite nonnegative combination of atoms $\Phi_{B,t,\tau_j^*}$, this completes the proof.
\end{proof}

\begin{remark}[Parameter scaling]
The bandwidth $B$ and heat scale $t$ depend on the compact $[-R, R]$ and accuracy $\varepsilon$. No uniform bound $\sup_K B(K)$ is required for density; the schedules $K \mapsto B(K)$, $K \mapsto t(K)$ may grow with $K$.
\end{remark}

%=============================================================================
\section{Lipschitz continuity of $Q$ (A2)}
%=============================================================================

\subsection{Finiteness of the prime sampler}

\begin{lemma}[Local finiteness]\label{lem:local-finite}
Fix $K > 0$. For $\Phi \in C_c([-K, K])$ even, the prime sum
\[
  \sum_{n \geq 2} w_Q(n)\, \Phi(\xi_n)
\]
is a finite sum: only $n \leq e^{2\pi K}$ contribute.
\end{lemma}

\begin{proof}
$\Phi(\xi_n) = 0$ for $|\xi_n| > K$, i.e., for $(\log n)/(2\pi) > K$, equivalently $n > e^{2\pi K}$.
\end{proof}

\subsection{The Lipschitz bound}

\begin{proof}[Proof of Theorem~\ref{thm:A2-Lip}]
Write $Q(\Phi_1) - Q(\Phi_2) = I_A + I_P$ where
\[
  I_A = \int_{-R}^R a^*(\xi)\, (\Phi_1 - \Phi_2)(\xi)\, d\xi, \qquad
  I_P = -\sum_{\xi_n \in [-R,R]} w_Q(n)\, (\Phi_1 - \Phi_2)(\xi_n).
\]

\textbf{Archimedean term.} Since $a^*$ is continuous on $[-R, R]$:
\[
  |I_A| \leq \int_{-R}^R |a^*(\xi)|\, |\Phi_1 - \Phi_2|(\xi)\, d\xi \leq \|a^*\|_{L^1(K)}\, \|\Phi_1 - \Phi_2\|_\infty.
\]

\textbf{Prime term.} By Lemma~\ref{lem:local-finite}, the sum is finite:
\[
  |I_P| \leq \sum_{\xi_n \in K} w_Q(n)\, |\Phi_1(\xi_n) - \Phi_2(\xi_n)| \leq \left( \sum_{\xi_n \in K} w_Q(n) \right) \|\Phi_1 - \Phi_2\|_\infty.
\]

Adding the bounds gives~\eqref{eq:Lip-bound} with $L_Q(K) = \|a^*\|_{L^1(K)} + \sum_{\xi_n \in K} w_Q(n)$.
\end{proof}

\subsection{Tail control for Fej\'er--heat windows}

\begin{lemma}[Gaussian tail bound]\label{lem:tail}
For $\Phi_{B,t}$ a Fej\'er--heat window and $N \geq e^2$:
\[
  \mathrm{Tail}(t; N) := \sum_{n > N} w_Q(n)\, \Phi_{B,t}(\xi_n) \leq \frac{1}{t}\, e^{-t(\log N)^2}.
\]
\end{lemma}

\begin{proof}
For the heat factor, $e^{-4\pi^2 t \xi_n^2} = e^{-t(\log n)^2}$. Thus
\[
  \sum_{n > N} w_Q(n)\, \Phi(\xi_n) \leq \sum_{n > N} \frac{2\log n}{\sqrt{n}}\, e^{-t(\log n)^2}.
\]
Via the integral comparison $\sum_{n > N} \leq \int_{N-1}^\infty$ and the substitution $y = \log x$:
\[
  \int_{\log N}^\infty y\, e^{-ty^2}\, dy = \frac{1}{2t}\, e^{-t(\log N)^2}.
\]
Multiplying by $2$ gives the stated bound.
\end{proof}

\begin{remark}[Leakage control]
When Fej\'er--heat windows are used on $[-K, K]$, the Gaussian factor produces exponentially small leakage outside the compact. This tail is absorbed into the (A2) continuity budget.
\end{remark}

%=============================================================================
\section{Discussion}
%=============================================================================

\subsection{Summary}

We have established three foundational results:
\begin{enumerate}
  \item \textbf{T0 (Normalization)}: The repository convention $\xi = \eta/(2\pi)$ matches Guinand--Weil with Jacobian factor $a^* = 2\pi a$.
  \item \textbf{A1' (Density)}: Fej\'er--heat atoms generate a cone dense in $C^+_{\mathrm{even}}([-K,K])$.
  \item \textbf{A2 (Lipschitz)}: $Q$ is Lipschitz with explicit constant $L_Q(K)$ on each compact.
\end{enumerate}

\subsection{Relation to Part~II}

The density (A1') and continuity (A2) results reduce the positivity problem $Q \geq 0$ on the full Weil cone to verifying $Q \geq 0$ on Fej\'er--heat atoms. In [Part~II], we establish:
\begin{itemize}
  \item \textbf{A3 (Toeplitz barrier)}: The Archimedean symbol $P_A(\theta)$ satisfies $\inf_\theta P_A(\theta) \geq c_* = 11/10$.
  \item \textbf{RKHS cap}: The prime operator $T_P$ satisfies $\|T_P\|_{\mathcal{H}_t \to \mathcal{H}_t} < c_*/4$.
\end{itemize}
Together, these imply $Q(\Phi) \geq (c_* - c_*/4)\|\Phi\|^2 > 0$ for atoms.

\subsection{Open questions}

\begin{enumerate}
  \item Can the Lipschitz constant $L_Q(K)$ be computed in closed form for small $K$?
  \item What is the optimal heat scale $t$ minimizing the total approximation error in (A1')?
  \item Are there other test cones (beyond Fej\'er--heat) with similar density and continuity properties?
\end{enumerate}

%=============================================================================
% Bibliography
%=============================================================================

\begin{thebibliography}{99}

\bibitem{Guinand1948}
A.~P. Guinand,
\textit{A summation formula in the theory of prime numbers},
Proc. London Math. Soc. (2) \textbf{50} (1948), 107--119.

\bibitem{Weil1952}
A.~Weil,
\textit{Sur les ``formules explicites'' de la th\'eorie des nombres premiers},
Comm. S\'em. Math. Univ. Lund (1952), 252--265.

\bibitem{SteinShakarchi2003}
E.~M. Stein and R.~Shakarchi,
\textit{Fourier Analysis: An Introduction},
Princeton University Press, 2003.

\bibitem{NISTDLMF}
NIST Digital Library of Mathematical Functions,
\url{https://dlmf.nist.gov}, Release 1.1.

\end{thebibliography}

\end{document}
